\documentclass[10pt]{mypackage}

%\usepackage{mlmodern}
%\usepackage{newpxtext,eulerpx,eucal}
%\renewcommand*{\mathbb}[1]{\varmathbb{#1}}
\usepackage[light]{kpfonts}
\usepackage{homework}
%\usepackage{notes}

%\usepackage[ backend=bibtex, style = alphabetic, sorting=ynt ]{biblatex}
%\addbibresource{  }

\usepackage{parskip}

\fancyhf{}
\fancyhead[R]{Avinash Iyer}
\fancyhead[L]{Algebraic Topology: Homework 20}
\fancyfoot[C]{\thepage}

\setcounter{secnumdepth}{0}

\begin{document}
\RaggedRight
\begin{problem}
  Complete the proof from class that if $f,g\colon X\rightarrow Y$ are homotopic maps of spaces, then the induced maps $f_{\#},g_{\#}\colon C_{\ast}\left( X \right)\rightarrow C_{\ast}\left( Y \right)$ are chain homotopic maps of chain complexes.
\end{problem}
\begin{solution}
  Recall that the prism operator $P_q\left( \sigma \right)$ is defined by its action on singular $k$-simplices $\sigma\colon \left[ e_0,\dots,e_n \right]\rightarrow X$, given by
  \begin{align*}
    P_q\left( \sigma \right) &= \sum_{i=0}^{q} \left( -1 \right)^{i} H\circ \left( \sigma\times \id \right)|_{\left[ v_0,\dots,v_i,w_i,\dots,w_q \right]}.
  \end{align*}
  We recall that $w_j = g\left( \sigma\left( e_j \right) \right)$ and $v_i = f\left( \sigma\left( e_i \right) \right)$. This allows us to rewrite the prism operator as
  \begin{align*}
    P_q\left( \sigma \right) &= \sum_{i=0}^{q} \left( -1 \right)^{i} \left[ f\left( \sigma\left( e_0 \right) \right),\dots,f\left( \sigma\left( e_i \right) \right),g\left( \sigma\left( e_i \right) \right),\dots,g\left( \sigma\left( e_q \right) \right) \right].
  \end{align*}
  To compute $\partial_{q+1}P_q\left( \sigma \right) + P_{q-1}\partial_{q}\left( \sigma \right)$, we do so in steps. We have
  \begin{align*}
    P_{q-1}\partial_q\left( \sigma \right) &= P_{q-1}\left( \sum_{i=0}^{q}\left( -1 \right)^{i}\left[ \sigma\left( e_0 \right),\dots,\widehat{\sigma\left( e_i \right)},\dots,\sigma\left( q \right) \right] \right)\\
                                           &= \sum_{i=0}^{q}\sum_{j < i}\left( -1 \right)^{i+j}\left[ f\left( \sigma\left( e_0 \right) \right),\dots, f\left( \sigma\left( e_j \right) \right),g\left( \sigma\left( e_j \right) \right),\dots,\widehat{g\left( \sigma\left( e_i \right) \right)},\dots,g\left( \sigma\left( e_q \right) \right) \right]\tag{A1}\\
                                           &+ \sum_{i=0}^{q} \left[ f\left( \sigma\left( e_0 \right) \right),\dots,f\left( \sigma\left( e_i \right) \right),g\left( \sigma\left( e_{i+1} \right) \right),\dots,g\left( \sigma\left( e_q \right) \right) \right]\tag{A2}\\
                                           &+ \sum_{i=0}^{q}\sum_{j> i}\left( -1 \right)^{i+j}\left[ f\left( \sigma\left( e_0 \right) \right),\dots,\widehat{f\left( \sigma\left( e_i \right) \right)},\dots,f\left( \sigma\left( e_j \right) \right),g\left( \sigma\left( e_j \right) \right),\dots,g\left( \sigma\left( e_q \right) \right) \right]\tag{A3}\\
    \partial_{q+1}P_q\left( \sigma \right) &= \partial_{q+1}\left( \sum_{i=0}^{q}\left( -1 \right)^{i} \left[ f\left( \sigma\left( e_0 \right) \right),\dots,f\left( \sigma\left( e_i \right) \right),g\left( \sigma\left( e_i \right) \right),\dots,g\left( \sigma\left( e_q \right) \right) \right] \right)\\
                                           &= \sum_{i=0}^{q}\left( -1 \right)^{i} \partial_{q+1} \left( \left[ f\left( \sigma\left( e_0 \right) \right),\dots,f\left( \sigma\left( e_i \right) \right),g\left( \sigma\left( e_i \right) \right),\dots,g\left( \sigma\left( e_q \right) \right) \right] \right)\\
                                           &= \sum_{i=0}^{q}\sum_{j < i} \left( -1 \right)^{i+j-1} \left[ f\left( \sigma\left( e_0 \right) \right),\dots,\widehat{f\left( \sigma\left( e_j \right) \right)},\dots,f\left( \sigma\left( e_i \right) \right),g\left( \sigma\left( e_i \right) \right),\dots,g\left( \sigma\left( e_q \right) \right) \right]\tag{B1}\\
                                           &+ \left( -1 \right)\sum_{i=0}^{q} \left[ f\left( \sigma\left( e_0 \right) \right),\dots,f\left( \sigma\left( e_{i-1} \right) \right),g\left( \sigma\left( e_i \right) \right),\dots,g\left( \sigma\left( e_q \right) \right) \right]\tag{B2}\\
                                           &+ \sum_{i=0}^{q}\sum_{j > i} \left( -1 \right)^{i+j-1} \left[ f\left( \sigma\left( e_0 \right) \right),\dots,f\left( \sigma\left( e_i \right) \right),g\left( \sigma\left( e_i \right) \right),\dots, \widehat{g\left( \sigma\left( e_j \right) \right)},\dots,g\left( \sigma\left( e_q \right) \right) \right].\tag{B3}
  \end{align*}
  Here, the extra factor of $-1$ in equations (B1) through (B3) comes from the fact that we are considering the boundary of $\left( q+1 \right)$-simplices.

  As an example for the cancellations that will be seen when we take $\partial_{q+1}P_q(\sigma) + P_{q-1}\partial_q(\sigma)$, observe that in (A1), we have that if $i = 1$ and $j = 0$, we have the simplex
  \begin{align*}
    -\left[ f\left( \sigma\left( e_0 \right) \right),g\left( \sigma\left( e_0 \right) \right),\widehat{g\left( \sigma\left( e_1 \right) \right)},\dots,g\left( \sigma\left( e_q \right) \right) \right],
  \end{align*}
  while in (B3), if $i = 0$ and $j = 1$, we have
  \begin{align*}
    \left[ f\left( \sigma\left( e_0 \right) \right),g\left( \sigma\left( e_0 \right) \right),\widehat{g\left( \sigma\left( e_1 \right) \right)},\dots,g\left( \sigma\left( e_q \right) \right) \right].
  \end{align*}
  Furthermore, the intermediate terms in both (A2) and (B2) cancel out to yield the final value 
  \begin{align*}
    \left[ f\left( \sigma\left( e_0 \right) \right),\dots,f\left( \sigma\left( e_q \right) \right) \right] - \left[ g\left( \sigma\left( e_0 \right) \right),\dots,g\left( \sigma\left( e_q \right) \right) \right],
  \end{align*}
  which is precisely $f(\sigma)-g(\sigma)$.
\end{solution}
\begin{problem}[Problem 2]
  Prove that chain homotopy defines an equivalence relation on the set of chain maps between fixed chain complexes $C_{\ast}$ and $D_{\ast}$.
\end{problem}
\begin{solution}
  Suppose $f_{\#},g_{\#},h_{\#}\colon C_{\ast}\rightarrow D_{\ast}$ are chain maps.

  First, we observe that for all singular $q$-simplices $\sigma$, we have $f(\sigma)-f(\sigma) = 0$, meaning that the chain map $P\colon C_{q}\rightarrow D_{q+1}$ given by $P(\sigma) = 0$ defines a chain homotopy between $f$ and itself. Therefore, chain homotopy is reflexive.

  For symmetry, if $P_q$ is a chain homotopy between $f_{q}$ and $g_q$, then $-P_q$ satisfies
  \begin{align*}
    \partial'\left( -P_q \right)(\sigma) + \left( -P_{q-1} \right)\partial(\sigma) &= - \left( \partial'P_q(\sigma) + P_{q-1}\partial(\sigma) \right)\\
                                                                                   &= -\left( f_q(\sigma)-g_q(\sigma) \right)\\
                                                                                   &= g_{q}(\sigma)-f_q(\sigma),
  \end{align*}
  meaning that $-P_q$ is a chain homotopy between $g_q$ and $f_q$.

  Finally, for transitivity, we let $P_q$ be a chain homotopy between $f_q$ and $g_q$, and $P_q'$ a chain homotopy between $g_q$ and $h_q$. These satisfy
  \begin{align*}
    \partial_{q+1}'P_q + P_{q-1}\partial_q &= f_q - g_q\\
    \partial_{q+1}'P_q' + P_{q-1}'\partial_q &= g_q - h_q,
  \end{align*}
  so since $P_q$ and $P_q'$ are homomorphisms of abelian groups, we may add them together, and observe that the map $P_q + P_q'$ is a map between $C_q$ and $D_{q+1}$ satisfying
  \begin{align*}
    \partial_{q+1}'\left( P_q + P_q' \right) + \left( P_{q-1} +P_{q-1}' \right)\partial_q &= \left( \partial_{q+1}'P_q + P_{q-1}\partial_q \right) + \left( \partial_{q+1}'P_q' + P_{q-1}'\partial_q \right)\\
                                                                                          &= \left( f_q - g_q \right) + \left( g_q - h_q \right)\\
                                                                                          &= f_q - h_q,
  \end{align*}
  so the relation ``is chain-homotopic to'' is an equivalence relation.
\end{solution}
\end{document}
